% Clase 20 — García Olvera Emiliano (actuadores)
\paragraph{Actuadores}
\textbf{Autor:} García Olvera Emiliano

Las órdenes llegan en un reporte de salida HID de 8 bytes (\texttt{uint8\_t}). Las duraciones y frecuencias viajan
como \texttt{uint16\_t} en \textit{little-endian} (máximo 65\,535) y se arman con \texttt{le16()}; el firmware acepta
de 1 a 5000 ms y de 100 a 5000 Hz (\texttt{actuadores.h}). El \textit{wrap} del PWM del buzzer también es de 16 bits:
\begin{align*}
\text{wrap} &= \frac{125\text{ MHz}}{64 \times f} - 1 \\
100\text{ Hz} &\rightarrow 19\,530 \\
220\text{ Hz} &\rightarrow 8876 \\
5000\text{ Hz} &\rightarrow 389
\end{align*}
y el peor caso (19\,530) cabe en 16 bits. El motor usa un PWM fijo de $125\text{ MHz} / (6.25 \times 1000) = 20\text{ kHz}$,
por encima de lo audible.

\texttt{volatile} está en \texttt{s\_alarma\_motor} y \texttt{s\_alarma\_buzzer} (\texttt{actuadores.c}, líneas 19 y 20):
las pone en 0 la alarma (una interrupción del temporizador) cuando apaga el actuador, y las lee el código normal
para cancelar la alarma anterior si llega otra orden. Sin \texttt{volatile}, el compilador podría usar un valor
viejo y cancelar una alarma que ya no existe.
